\chapter{Комбинаторика}

\section{Основные формулы}
Выбор $k$ элементов из $n$ различных:
\begin{align*}
	P_n &= n! && \text{перестановки}\\
	A_n^k &= \frac{n!}{(n-k)!} && \text{размещения}\\
	\bar A_n^k &= n^k && \text{размещения с повторениями}\\
	C_n^k = \binom{n}{k} &= \frac{n!}{k!\,(n-k)!} && \text{сочетания}
\end{align*}

\subsection{Свойства биномиальных коэффициентов}
Для $0 \le k \le n$:
\begin{align*}
	\binom{n}{k} &= \binom{n}{n-k} = \frac{n}{k}\binom{n-1}{k-1}\\
	\binom{n}{k} &= \binom{n-1}{k-1} + \binom{n-1}{k} \quad \text{(Паскаль)}\\
	\sum_{k=0}^{n} \binom{n}{k} &= 2^n, \qquad \sum_{k=0}^{n} (-1)^k \binom{n}{k} = 0 \quad (n > 0)\\
	\sum_{k=0}^{n} \binom{n}{k}^2 &= \binom{2n}{n}, \qquad \sum_{i=k}^{n} \binom{i}{k} = \binom{n+1}{k+1}\\
	\sum_{i=0}^{k} \binom{m}{i}\binom{n}{k-i} &= \binom{m+n}{k} \quad \text{(Вандермонд)}\\
	(x + y)^n &= \sum_{k=0}^{n} \binom{n}{k} x^k y^{n-k}
\end{align*}

\textbf{Теорема Люка.} Для простого $p$, где $n_i, k_i$ --- цифры $n$ и $k$ в системе счисления по основанию $p$:
\[ \binom{n}{k} \equiv \prod_i \binom{n_i}{k_i} \pmod p. \]

\section{Сочетания с повторениями}
Число способов выбрать $k$ элементов из $n$ типов (порядок не важен, повторы разрешены), то есть число мультимножеств размера $k$:
\[ \bar C_n^k = \binom{n+k-1}{k} = \binom{n+k-1}{n-1}. \]
\textbf{Шары и перегородки.} То же самое --- число решений уравнения $x_1 + \dots + x_n = k$ в целых $x_i \ge 0$ ($k$ шаров и $n-1$ перегородка).
\begin{itemize}
	\item $x_i \ge 1$: \ $\binom{k-1}{n-1}$.
	\item $x_i \ge a_i$: замена $y_i = x_i - a_i$ даёт $\binom{k - \sum a_i + n - 1}{n-1}$.
	\item $x_1 + \dots + x_n \le k$, $x_i \ge 0$: \ $\binom{n+k}{n}$ (фиктивная переменная $x_{n+1}$).
	\item $0 \le x_i \le b$: по формуле включений-исключений
	\[ \sum_{j \ge 0} (-1)^j \binom{n}{j} \binom{k - j(b+1) + n - 1}{n-1}. \]
\end{itemize}

\section{Мультиномиальные коэффициенты}
Для $k_1 + \dots + k_m = n$:
\[ \binom{n}{k_1, k_2, \dots, k_m} = \frac{n!}{k_1!\,k_2!\cdots k_m!}. \]
\begin{itemize}
	\item Число различных слов из $k_1$ букв первого вида, \dots, $k_m$ букв $m$-го вида (перестановки с повторениями).
	\item Число способов разложить $n$ различных предметов по $m$ подписанным ящикам, в $i$-й ящик --- ровно $k_i$ предметов.
	\item Если ящики неразличимы и одного размера $k$ ($n = mk$), делим на $m!$: \ $\frac{n!}{(k!)^m\, m!}$.
\end{itemize}
Через биномиальные коэффициенты (удобно считать по модулю):
\[ \binom{n}{k_1, \dots, k_m} = \binom{k_1}{k_1}\binom{k_1+k_2}{k_2}\cdots\binom{k_1+\dots+k_m}{k_m}. \]
Рекуррентность:
\[ \binom{n}{k_1, \dots, k_m} = \sum_{i=1}^{m} \binom{n-1}{k_1, \dots, k_i - 1, \dots, k_m}. \]
\textbf{Полиномиальная теорема:}
\[ (x_1 + \dots + x_m)^n = \sum_{k_1 + \dots + k_m = n} \binom{n}{k_1, \dots, k_m} x_1^{k_1} \cdots x_m^{k_m}. \]
Сумма всех мультиномиальных коэффициентов равна $m^n$, а число слагаемых равно $\binom{n+m-1}{m-1}$.

\section{Перестановки}
\textbf{Беспорядки} (перестановки без неподвижных точек):
\[ D_n = (n-1)(D_{n-1} + D_{n-2}) = n! \sum_{k=0}^{n} \frac{(-1)^k}{k!} = \left\lfloor \frac{n!}{e} + \frac{1}{2} \right\rfloor. \]
$D_n$: 1, 0, 1, 2, 9, 44, 265, 1854, 14833.

\textbf{Числа Стирлинга I рода} $c(n, k)$ --- перестановки $n$ элементов ровно с $k$ циклами:
\[ c(n, k) = c(n-1, k-1) + (n-1)\, c(n-1, k), \qquad \sum_k c(n, k) = n!. \]

\section{Разбиения}
\textbf{Числа Стирлинга II рода} $S(n, k)$ --- разбиения $n$ различных элементов на $k$ непустых неупорядоченных групп:
\[ S(n, k) = k\, S(n-1, k) + S(n-1, k-1) = \frac{1}{k!} \sum_{j=0}^{k} (-1)^j \binom{k}{j} (k-j)^n. \]
Число сюръекций $n \to k$ равно $k!\, S(n, k)$.

\textbf{Числа Белла} --- все разбиения множества из $n$ элементов:
\[ B_n = \sum_{k} S(n, k), \qquad B_{n+1} = \sum_{k=0}^{n} \binom{n}{k} B_k. \]
$B_n$: 1, 1, 2, 5, 15, 52, 203, 877, 4140, 21147.

\section{Числа Каталана}
\[ C_n = \frac{1}{n+1}\binom{2n}{n} = \binom{2n}{n} - \binom{2n}{n+1}, \qquad C_{n+1} = \sum_{i=0}^{n} C_i C_{n-i}. \]
$C_n$: 1, 1, 2, 5, 14, 42, 132, 429, 1430, 4862, 16796.

$C_n$ --- это число правильных скобочных последовательностей из $n$ пар, путей по сетке $n \times n$ не выше диагонали, бинарных деревьев с $n$ вершинами, триангуляций $(n+2)$-угольника.

\section{Подсчёт с симметриями}
\textbf{Формула включений-исключений:}
\[ \Big|\bigcup_{i=1}^{n} A_i\Big| = \sum_{\varnothing \ne S \subseteq \{1..n\}} (-1)^{|S|+1} \Big|\bigcap_{i \in S} A_i\Big|. \]
\textbf{Лемма Бёрнсайда:} число классов эквивалентности относительно группы $G$
\[ \frac{1}{|G|} \sum_{g \in G} |X^g|, \]
где $X^g$ --- множество объектов, неподвижных при действии $g$. Ожерелья из $n$ бусин $k$ цветов с точностью до поворота:
\[ \frac{1}{n} \sum_{d \mid n} \varphi(d)\, k^{n/d}. \]
